\documentclass[11pt,a4paper]{article}
\usepackage[T1]{fontenc}
\usepackage{lmodern}
\usepackage[margin=25mm]{geometry}
\usepackage{amsmath,amssymb,amsthm,mathtools,booktabs}
\usepackage[numbers,sort&compress]{natbib}
\usepackage{microtype}
\usepackage[colorlinks=true,linkcolor=blue,citecolor=blue,urlcolor=blue]{hyperref}
\hypersetup{pdftitle={An eighth-order inverse-response filter in a ten-site quantum walk},
 pdfauthor={Xiaojian Zeng},pdfsubject={Fixed-chain postselected matrix inversion},
 pdfkeywords={quantum walk, matrix inversion, interval arithmetic, postselection}}
\usepackage{url}
\newcommand{\ii}{\mathrm{i}}
\newcommand{\dd}{\mathrm{d}}
\newcommand{\ket}[1]{\lvert #1\rangle}
\newcommand{\bra}[1]{\langle #1\rvert}
\newcommand{\norm}[1]{\lVert #1\rVert}
\newcommand{\Res}{\operatorname{Res}}
\newtheorem{theorem}{Theorem}
\newtheorem{lemma}{Lemma}
\numberwithin{equation}{section}
\title{An eighth-order inverse-response filter in a ten-site quantum walk}
\author{Xiaojian Zeng\\[3pt]
\small PBC School of Finance, Tsinghua University, 100190, China\\
\small\href{mailto:xiaojianzeng@foxmail.com}{xiaojianzeng@foxmail.com}}
\date{}
\begin{document}
\maketitle
\begin{abstract}
We construct a fixed, positive-coupling ten-site quantum walk with uniform
eighth-order absolute inverse-response error. An interval contraction certifies
simultaneous cancellation of the cubic, quintic and septic low-energy coefficients
at dimensionless time $2\pi$. A contour argument and a separate high-energy
spectral-weight estimate prove an explicit bound for $0<x\le1/16$, retaining the
rapidly oscillating contribution. The corresponding postselected inverse has
seventh-order conditional-state error under an explicit spectral-domain condition.
Matched propagation examples report success probability and Hamiltonian norm cost.
This resolves the four-extension existence question within the autonomous
quantum-walk inversion model; the general chain-length conjecture remains open.
\end{abstract}
\noindent\textbf{Keywords:} continuous-time quantum walk; matrix inversion;
spectral filter; interval arithmetic; postselection

\section{Introduction}
Weakly coupled quantum walks provide inverse-matrix responses through one
evolution under a fixed Hamiltonian. In their preprint
\emph{Matrix Inversion by Quantum Walk}
(\href{https://arxiv.org/abs/2508.06611v1}{arXiv:2508.06611v1}, 2025),
Kay and Tamon formulate an existence conjecture for chains with $R$ extensions
on either side of the matrix block. Their Sec.~II.A, Table~I and Appendix~A
give a ten-site, $R=4$ design with absolute response error of order $x^7$
under a fixed low-energy spectrum. The desired unrestricted order is $x^8$.
Their odd low-energy expansion and suppression of the high-energy endpoint
weight supply the mechanism underlying the present construction.

We free one finite spectral eigenvalue and simultaneously cancel the cubic,
quintic and septic coefficients. An exact root specified by an interval
contraction fixes four positive couplings. We then bound the entire amplitude
uniformly, including the high-energy oscillation. An ordinary Taylor expansion
of the full Hamiltonian at $x=0$ would miss this contribution and would not
establish the claimed error order.

Related weak-coupling transfer results are developed in Kay and Tamon's
\emph{The strength of weak coupling}
(\href{https://arxiv.org/abs/2512.08141v1}{arXiv:2512.08141v1}, 2025),
Theorem~4.1 and Secs.~4.2 and 6. Quantitative transfer estimates on graphs with
an involution are also studied by Lippner and Shi~\cite{LippnerShi2024}.
These transfer settings differ from the fixed-time inverse filter considered
here. Kang and Su~\cite{KangSu2026}, Sec.~4.4, Proposition~19 and
Corollary~20 of their current preprint version, construct inverses through
Hamiltonian singular-value transformation using repeated queries. That
resource setting is distinct from one fixed autonomous chain. Our comparison
concerns the nearest checked fixed-chain design; it makes no exhaustive
priority, general-$R$, or complete algorithmic speedup claim.

\section{Model and uniform response result}
In an orthonormal single-excitation site basis, let
\begin{equation}
 T_4(x;J)=\sum_{j=1}^{9}w_j
 \bigl(\ket{j}\bra{j+1}+\ket{j+1}\bra{j}\bigr),
 \qquad x>0,
 \label{eq:chain}
\end{equation}
where the open chain has zero diagonal and
\begin{equation}
 (w_1,\ldots,w_9)=(J_1,J_2,J_3,J_4,x^{-1},J_4,J_3,J_2,J_1).
\end{equation}
Its full endpoint amplitude is
\begin{equation}
 a_4(x;J,\tau)=\bra{10}e^{-\ii\tau T_4(x;J)}\ket{1}.
 \label{eq:amplitude}
\end{equation}
All four couplings and the dimensionless time are fixed independently of $x$.

\begin{theorem}\label{thm:main}
There are positive fixed couplings $J_1,\ldots,J_4$ and a real
$\ell\in(-1.078612,-1.078611)$ such that, at $\tau=2\pi$,
\begin{equation}
 \bigl|a_4(x;J,2\pi)-\ii\ell x\bigr|
 \le 131072002\,x^8,\qquad 0<x\le\frac1{16}.
 \label{eq:uniform}
\end{equation}
The couplings are specified exactly by the unique cancellation root in
Sec.~\ref{sec:root}. For orientation, they are approximately
\begin{equation}
 (J_1,J_2,J_3,J_4)\simeq
 (0.6803742065,\ 0.9435755016,\ 1.1381740446,\ 0.8283449422).
 \label{eq:rounded}
\end{equation}
The displayed rounded values are not a substitute for the exact root.
\end{theorem}

The numerical constant in Eq.~\eqref{eq:uniform} is deliberately conservative.
The conclusion is an existence statement in the finite Hermitian model,
with a computer-assisted root certificate and analytic remainder estimates.
It covers the $R=4$ case, without a construction for arbitrary $R$.

\section{Reflection sectors and cancellation equations}
\subsection{Exact resolvent representation}
Reflection about the central bond gives sectors spanned by
$(\ket{j}\pm\ket{11-j})/\sqrt{2}$, $j=1,\ldots,5$.
The two blocks $B_\pm(x)$ have off-diagonals $J_1,\ldots,J_4$,
and their only nonzero diagonal entry is the last, $\pm x^{-1}$.
For $D=\operatorname{diag}(1,-1,1,-1,1)$,
$B_-(x)=-D B_+(x)D$. Therefore
\begin{align}
 a_4(x;J,\tau)
 &=\frac12\left(\bra{1}e^{-\ii\tau B_+(x)}\ket{1}
             -\bra{1}e^{-\ii\tau B_-(x)}\ket{1}\right)\nonumber\\
 &=-\ii\sum_{r=1}^{5}W_r(x)\sin\bigl(\tau E_r(x)\bigr),
 \label{eq:sectors}
\end{align}
where $E_r$ are the eigenvalues of $B_+$ and $W_r$ their first-site
spectral weights. In particular, no extra factor of two enters the endpoint
response.

Write $a=J_1^2$, $b=J_2^2$, $c=J_3^2$, $d=J_4^2$ and set
\begin{align}
 p(z)&=z^4-(a+b+c)z^2+ac,
 &q(z)&=zp(z)-dz(z^2-a-b),\nonumber\\
 n_0(z)&=z^3-(b+c)z,
 &n_1(z)&=zn_0(z)-d(z^2-b).
 \label{eq:polynomials}
\end{align}
Continuant recursion for the determinant and first-site minor gives
\begin{equation}
 \det(zI-B_+)=-\frac{p(z)-xq(z)}{x},\qquad
 G_{11}(z,x)=\frac{n_0(z)-xn_1(z)}{p(z)-xq(z)}.
 \label{eq:resolvent}
\end{equation}
The polynomial identity $n_0q-n_1p=-abcd$ yields
\begin{equation}
 G_{11}(z,x)=\frac{n_0(z)}{p(z)}
 -abcd\sum_{n\ge1}x^n\frac{q(z)^{n-1}}{p(z)^{n+1}},
 \qquad |xq/p|<1.
 \label{eq:geometric}
\end{equation}

Let $a_{\mathrm{low}}=\ii f(x)$ denote the contribution from the four
finite eigenvalues inside a fixed contour. The paired zero-$x$ spectrum
makes $f_0=0$. Contour integration of Eq.~\eqref{eq:geometric} gives
\begin{equation}
 f(x)=\sum_{n\ge1}f_nx^n,\qquad
 f_n=abcd\,S_n,\qquad
 S_n=\sum_{p(\rho)=0}\Res_{z=\rho}
       \frac{\sin(\tau z)q(z)^{n-1}}{p(z)^{n+1}}.
 \label{eq:coefficients}
\end{equation}
Since $p$ is even and $q$ is odd, the integrand is even for even $n$,
and opposite-root residues cancel. Thus every even coefficient vanishes.
Only the low-energy part is expanded here.

\subsection{An exactly specified positive root}\label{sec:root}
Fix $u=1/2$ and $\tau=2\pi$. Parametrize
\begin{align}
 p(z)&=(z^2-u^2)(z^2-v^2),\nonumber\\
 a&=\frac{u^2v^2}{c},\qquad
 b=u^2+v^2-c-\frac{u^2v^2}{c},\qquad
 s=a+b=u^2+v^2-c.
 \label{eq:parametrization}
\end{align}
The domain $u^2<c<v^2$, $v>u$, $d>0$ makes all four squared
couplings positive: $bc=(c-u^2)(v^2-c)>0$.
The three equations are
\begin{equation}
 F(v,c,d)=\bigl(S_3,S_5,S_7\bigr)=0.
 \label{eq:cancellation}
\end{equation}
They are real analytic on the box used below. Each residue is a finite
polynomial-jet calculation. More explicitly, for $\rho\in\{u,v\}$ and
the other positive root $\omega$,
\begin{align}
 p(\rho+y)&=y h_\rho(y),\nonumber\\
 h_\rho(y)&=2\rho(\rho^2-\omega^2)
          +(5\rho^2-\omega^2)y+4\rho y^2+y^3,\label{eq:jet}\\
 S_n&=2\sum_{\rho=u,v}[y^n]
 \left\{\sin\bigl(2\pi(\rho+y)\bigr)
 q(\rho+y)^{n-1}h_\rho(y)^{-n-1}\right\},\qquad n\ \text{odd}.
 \label{eq:jetresidue}
\end{align}
This makes the cancellation equations reproducible without fitting a
singular full propagator.

Let $X$ be the closed infinity-norm box of radius $r=10^{-40}$ about
the exact terminating decimal vector $y_0=(v_0,c_0,d_0)$:
\begin{equation}
\begin{aligned}
v_0&=\texttt{1.5487685249240345145585943938889954395158348752808834652101731657}\\
c_0&=\texttt{1.2954401557316677104527788216689492776181313419030944437678424473}\\
d_0&=\texttt{0.68615534328287748264002673279247284550721282248356859067322025627}
\end{aligned}
\label{eq:center}
\end{equation}

Using a fixed exact decimal preconditioner $Y$ (Appendix~\ref{app:certificate}),
define $\Phi(y)=y-YF(y)$. Elementary directed interval arithmetic at
80 decimal digits, with an analytically differentiated jet recurrence,
encloses the complete Jacobian on $X$ and gives
\begin{align}
 \sup_{y\in X}\norm{I-YDF(y)}_\infty&<8.0\times10^{-36},\nonumber\\
 \norm{YF(y_0)}_\infty&<4.28\times10^{-65},\nonumber\\
 \sup_{y\in X}\norm{\Phi(y)-y_0}_\infty&<4.28\times10^{-65}<r.
 \label{eq:contraction}
\end{align}
The last line includes the Jacobian contribution times $r$, rather than
only the center displacement. Thus $\Phi$ maps $X$ strictly into itself
and is a contraction. Banach's fixed-point theorem gives a unique fixed
point. The first inequality also makes $YDF(y)$ invertible; hence $Y$
is invertible and the fixed point is an exact zero of $F$.
Conversely every zero is a fixed point, establishing uniqueness within $X$.

At this zero, define $J=(\sqrt{a},\sqrt{b},\sqrt{c},\sqrt{d})$ using
Eq.~\eqref{eq:parametrization}, and $\ell=abcd\,S_1$.
The same interval calculation gives positivity and
\begin{equation}
 -1.078612<\ell<-1.078611.
 \label{eq:linear}
\end{equation}
Consequently $f_3=f_5=f_7=0$ exactly, and
$f(x)=\ell x+O(x^9)$. These cancellations hold at the enclosed zero;
the terminating decimal center has small nonzero residuals.

The supplementary certificate retains all exact binary interval endpoints.
The enclosure uses only elementary arithmetic, sine, cosine and $\pi$
from \texttt{mpmath.iv} 1.3.0, with $\sin\pi=0$ and $\cos\pi=-1$
inserted exactly at $u=1/2$. Numerical root finding and inversion propose
the center and $Y$; their numerical accuracy is not an existence premise.
This is a computer-assisted proof relying on the interval implementation,
without proof-assistant formalization. A separate calculation from implicit
finite-eigenvalue series and cofactor weights reproduces the coefficients
through order seven; that numerical agreement supports the implementation
and does not replace the contraction inequalities.

\section{Uniform control of the complete amplitude}\label{sec:bound}
\subsection{The analytic low-energy remainder}
The root box implies the coarse rational bounds
\begin{equation}
 \frac32<v<\frac{31}{20},\quad
 1.2954<c<1.2955,\quad 0<d<0.687,\quad
 0<s<1.36,\quad 0<abcd<1.
 \label{eq:outer}
\end{equation}
On the circle $|z|=2$,
\begin{align}
 |p(z)|&\ge(4-u^2)(4-v^2)>5,\nonumber\\
 \left|\frac{q(z)}{p(z)}\right|
 &\le 2+\frac{2d(4+s)}{5}<4.
 \label{eq:contour}
\end{align}
For $|x|\le1/16$, $|xq/p|<1/4$. Rouch\'e's theorem counts exactly
four zeros of $p-xq$ inside this circle. The contour integral is analytic
in $x$ and Eq.~\eqref{eq:geometric} converges uniformly there, so
coefficient extraction under the integral is justified.

Since $|\sin(2\pi z)|\le e^{4\pi}<300000$, the circle length divided
by $2\pi$ is 2, and $abcd<1$, Eq.~\eqref{eq:coefficients} gives
\begin{equation}
 |f_n|\le 24000\,4^{n-1},\qquad n\ge1.
 \label{eq:coefficientbound}
\end{equation}
The elementary exponential bound can itself be checked rationally: use
$\pi<22/7$, sum $\sum_{k=0}^{100}(88/7)^k/k!$, and bound the
positive tail by its first term times $1/(1-(88/7)/102)$.
The resulting upper bound is less than 300000.
After the odd cancellations and even parity zeros, the remainder starts
at order nine. For real $0<x\le1/16$,
\begin{align}
 |f(x)-\ell x|
 &\le\frac{24000\,4^8x^9}{1-4x}\nonumber\\
 &\le2097152000\,x^9
 \le131072000\,x^8.
 \label{eq:lowbound}
\end{align}
Including the zero even coefficients in the geometric majorant only
enlarges this bound.

\subsection{The oscillatory high-energy branch}
For real positive $x$, $B_+(x)$ is Hermitian. Its largest eigenvalue
$\Lambda$ satisfies $\Lambda\ge1/x$ by the last-site Rayleigh quotient.
Since $1/x\ge16$, this is the fifth eigenvalue outside the contour.
Let $\eta$ be the last component of its normalized eigenvector.
Solving the first four component equations gives
\begin{equation}
 \bra{1}\psi_{\mathrm{high}}\rangle
 =\eta\frac{J_1J_2J_3J_4}{p(\Lambda)}.
 \label{eq:highcomponent}
\end{equation}
The numerator follows from the cross-chain cofactor of the four-site
principal block. With $|\eta|\le1$, its first-site weight obeys
\begin{align}
 W_{\mathrm{high}}
 &\le\frac{abcd}{p(\Lambda)^2}\nonumber\\
 &\le\frac{x^8}{\bigl[(1-u^2x^2)(1-v^2x^2)\bigr]^2}
 \le2x^8.
 \label{eq:highbound}
\end{align}
The last inequality uses $u^2<v^2<3$, $x^2\le1/256$ and
$(253/256)^4>1/2$.
Equation~\eqref{eq:sectors} identifies its full endpoint contribution as
\begin{equation}
 a_{\mathrm{high}}(x)=-\ii W_{\mathrm{high}}(x)
                      \sin\bigl(2\pi\Lambda(x)\bigr).
\end{equation}
Thus $|a_{\mathrm{high}}|\le2x^8$ for every $x$ in the specified
interval. No expansion or averaging of this phase has been used.
Adding Eq.~\eqref{eq:lowbound} proves Theorem~\ref{thm:main}.

For the specified witness, perturbation theory for the isolated eigenvalue
near 1 of $xB_+(x)$ gives $\Lambda=x^{-1}+dx+O(x^3)$ and
$W_{\mathrm{high}}=abcd\,x^8(1+O(x^2))$. Therefore
\begin{equation}
 a_4(x)-\ii\ell x
 =-\ii abcd\,x^8\sin\bigl(2\pi/x+O(x)\bigr)+O(x^9).
 \label{eq:asymptotic}
\end{equation}
Sequences with opposite sine maxima show
\begin{equation}
 \limsup_{x\to0^+}\frac{|a_4(x)-\ii\ell x|}{x^8}
 =abcd\simeq0.3663437609.
\end{equation}
This explains the attained order for this witness and the oscillatory
error envelope. It is not a general-chain optimality theorem.

\section{Postselected inverse, probability and cost}
Let $A=U\Sigma V^\dagger$ be square and invertible, with dimensionless
singular values in $[1/\kappa,1]$, $\kappa\ge1$, and $\norm{A}=1$.
In units $\hbar=1$, introduce a common energy scale $E_0$. A ten-block
Hamiltonian has the four couplings $\gamma J_j E_0$ on each side,
central block $E_0A$ and its adjoint, and zero block diagonal.
For each singular channel,
\begin{equation}
 H_\lambda/E_0=\gamma T_4(\gamma/\lambda;J),\qquad
 t=\frac{2\pi}{\gamma E_0},\qquad
 0<\gamma\le\frac{1}{16\kappa}.
 \label{eq:units}
\end{equation}
Thus $x=\gamma/\lambda$ remains inside the theorem's domain
simultaneously for the entire input spectrum. Sites before the central
block carry left singular vectors and sites after it carry right singular
vectors. Endpoint detection gives the Kraus block
\begin{equation}
 K=V\operatorname{diag}\bigl(a_4(\gamma/\lambda_j)\bigr)U^\dagger,
 \qquad K_0=\ii\ell\gamma A^{-1}.
\end{equation}
With $C=131072002$,
\begin{equation}
 \norm{K-K_0}\le\Delta,\qquad
 \Delta=C(\gamma\kappa)^8.
 \label{eq:krauserror}
\end{equation}
For normalized $b$, define $n=\norm{K_0b}=|\ell|\gamma\norm{A^{-1}b}$.
Triangle inequalities and the contraction property of $K$ give
\begin{equation}
 \max(0,n-\Delta)^2\le p_{\mathrm{succ}}=\norm{Kb}^2
 \le\min\bigl(1,(n+\Delta)^2\bigr).
 \label{eq:success}
\end{equation}
If $\Delta<n$, the successful branch is nonzero and
\begin{equation}
 \norm{\frac{Kb}{\norm{Kb}}-\frac{K_0b}{n}}
 \le\frac{2\Delta}{n}
 \le\frac{2C}{|\ell|}\gamma^7\kappa^8.
 \label{eq:conditional}
\end{equation}
For the first inequality, add and subtract $Kb/n$ and use
$|\norm{Kb}-n|\le\Delta$. The second uses
$\norm{A^{-1}b}\ge1$. The conditional reference includes the global
phase $\ii\ell/|\ell|$. The error power loses one order on normalization:
the absolute response order is eight, whereas this conditional-state bound
is seventh order for fixed $\kappa$ and input. Rectangular inverses require
a specified singular support and are not claimed here.

The complete evolution is unitary. The selected map
$\rho\mapsto K\rho K^\dagger$ is completely positive and trace
nonincreasing, with probability in Eq.~\eqref{eq:success}. The construction
does not yield deterministic inversion. Couplings stay of order one and
$\tau=2\pi$ matches the source design. A row-sum estimate gives the coarse
dimensionless cost bound
\begin{equation}
 \norm{H}t\le2\pi\left(\gamma^{-1}+2\max_jJ_j\right).
 \label{eq:cost}
\end{equation}
Unamplified attempts have expected repetition count $1/p_{\mathrm{succ}}$.
State preparation, Hamiltonian simulation, amplification, coherence and
control precision remain additional resources.

\section{Propagation comparison and limitations}
The fixed-spectrum reference keeps $u=1/2$, $v=3/2$ and numerically
cancels $S_3,S_5$ at high precision. Its couplings reproduce the rounded
$R=4$ values in Kay and Tamon's Table~I. Its response coefficient is about
$-0.9430116913$, and its septic coefficient is about $-0.2639054062$.
The new witness frees $v\simeq1.5487685249$ and cancels $S_7$ as well.
The largest coupling increases from about 0.9926 to 1.1382, while the
linear response magnitude increases to about 1.0786. Neither time nor
couplings diverge to obtain the higher order.

Table~\ref{tab:comparison} compares complete ten-site propagations at
$x=4/401$, where $2\pi/x$ is a sine maximum. These are scalar channel
diagnostics with identical $x$ and $\tau$. The prior reference is recomputed
at high precision rather than taken from six-digit printed couplings.
It is a precision-stabilized numerical comparison, not a separate
interval-certified prior theorem.
\begin{table}[htbp]
\centering
\small
\begin{tabular}{lccc}
\toprule
Design & $|a_4-\mathrm{i}\ell x|$ & $|a_4|^2$ & $\|T_4\|\tau$ \\
\midrule
Fixed-spectrum reference & $2.61870\times10^{-15}$ & $8.84841\times10^{-5}$ & $629.92122$ \\
Certified-root design & $3.27585\times10^{-17}$ & $1.15761\times10^{-4}$ & $629.93233$ \\
\bottomrule
\end{tabular}
\caption{Complete scalar-channel propagation at $x=4/401$ and $\tau=2\pi$.
The fixed-spectrum reference is reconstructed from the $R=4$ design in
Kay and Tamon, \emph{Matrix Inversion by Quantum Walk}
(arXiv:2508.06611v1), Table~I. Each error is measured relative to that
design's own linear coefficient. Displayed digits are numerical diagnostics;
the uniform bound follows from the root certificate and analytic estimates.}
\label{tab:comparison}
\end{table}


The calculations compare a direct ten-site representation with the
five-site spectral sum and repeat at 60 and 90 decimal digits.
Across six tested $x=4/(4k+1)$ values, $k=25,50,100,200,500,1000$,
the representation discrepancies are below $10^{-89}$ and the precision
differences below $7\times10^{-61}$. Coupling derivatives and Duhamel's
inequality bound the response uncertainty from the $10^{-40}$ parameter
box by $40\times10^{-40}$ over these propagations. These diagnostics
support the calculations; the interval root and uniform inequalities in
Sec.~\ref{sec:bound}, rather than the samples, prove Eq.~\eqref{eq:uniform}.
The leading scalar success-probability ratio is approximately
$(1.0786116373/0.9430116913)^2\simeq1.31$.

The sample shows substantially smaller response error at nearly equal
$\norm{H}t$. It does not establish a globally optimized filter, a complete
quantum algorithm advantage, or robustness to timing, coupling and hardware
errors. The uniform constant is loose; practical design would require
tighter constants and an implementation-error budget. Only the exact
$R=4$ existence slice is resolved. Existence and efficient discovery for
general $R$ remain outside the result. The nearest-primary-source
comparison is bounded and does not establish exhaustive priority.

\section*{Computational methods and AI disclosure}
The interval calculation uses \texttt{mpmath} 1.3.0. Symbolic determinant
and cofactor checks use \texttt{SymPy} 1.14.0; the numerical replay uses
arbitrary-precision Hermitian spectral propagation. Supplementary code
and exact certificate data specify the recurrences, precision settings,
parameter box and decimal preconditioner.
OpenAI Codex language-model tools were used for literature screening,
derivation and root-search assistance, verification-code generation,
computational interpretation and manuscript drafting. A separate AI
checking stage used a different coefficient representation; it is not
external peer review. This disclosure covers generative scientific work
as well as writing. Human scientific review and authorship confirmation
are pending for this local draft; completed human verification is not
asserted.

\section*{Statements and declarations}
\noindent\textbf{Data and code availability.}
The accompanying local supplement contains the root certificate,
propagation comparison data and scripts, symbolic checks, environment
specification and regeneration instructions. It contains no experimental
dataset. Public archival deposition has not occurred.\\[3pt]
\textbf{Funding and competing interests.}
These declarations await confirmation by the named author before any
submission.\\[3pt]
\textbf{Author information and contributions.}
The displayed author information is supplied for this draft and requires
confirmation. The named author's scientific review, contribution statement
and approval of the final text remain pending.\\[3pt]
\textbf{Ethics approval and consent.}
Not applicable to this finite-dimensional theoretical and computational study.

\appendix
\section{Exact interval certificate inputs}\label{app:certificate}
The fixed exact decimal preconditioner used in Eq.~\eqref{eq:contraction}
is listed row by row below. Each entry is a terminating decimal, interpreted
exactly before outward rounding. The entries and center were proposed
numerically; the strict interval inequalities provide the certificate.
\begingroup\small
\begin{verbatim}
Y11 = 0.061207121155273575213890779429074537437867825070659334072164057528
Y12 = -0.14614944776142125586130034028698605531319229338275959795286583887
Y13 = 0.033388744469063696323848108611820496559261062311112152234561659839
Y21 = 0.095988993500464343322236324835438325359250620898599511373109296823
Y22 = -0.48302764580534128969674320644336240119688446294722354899867335999
Y23 = 0.2065620366578123363538548095599313108560346650251941518385041284
Y31 = 0.34065699313566067153924085606325995484979187065078198062875093929
Y32 = -0.25729234664620376978066335496729766930071543693004756949961486848
Y33 = 0.14478866502255163895363617780142713794615098484178861690487505011
\end{verbatim}
\endgroup

The stored certificate supplies the full interval Jacobian and the exact
binary endpoints of $I-YDF(X)$, center corrections, image bounds,
coupling-positivity intervals and linear response. The replay instructions
evaluate the finite jets in Eq.~\eqref{eq:jetresidue} and differentiate
them analytically. A second representation solves the four implicit
eigenvalue series of $p(E)-xq(E)=0$, reconstructs first-site weights as
$(n_0(E)-xn_1(E))/(p'(E)-xq'(E))$, and composes their sine series.
This separates implementation checks from the analytic treatment of the
high-energy mode.

\bibliographystyle{unsrtnat}
\bibliography{references}
\end{document}
